\paragraph{Canceling one excitation is not stopping the Q-ball.}
The source-projection identity also supplies a simple linear control prediction.
Let $d=e^{-i\rho t}c$, initially zero, and apply two identical physical pulses
of the preceding family, with the same local rotating-frame waveform and
strength $\epsilon$, separated by $\Delta\geq1$.
Their total projected drive is
\begin{equation}
 I_{\mathrm{pair}}=\epsilon L(\delta)
                 \bigl(1+e^{-i\rho\Delta}\bigr).
 \label{env:counterpulse}
\end{equation}
Consequently $\Delta=\pi/\rho>1$ cancels the mode projection generated by the
first pulse, even for a common local detuning $\delta$. Both sideband terms
in $L$ are retained. This construction needs identical pulses in the
rotating frame: in laboratory variables their sources obey
$J_2(t)=e^{i\omega\Delta}J_1(t-\Delta)$.
A sign-reversed second pulse at that same half-period instead doubles the
rotating projected drive. For $|\delta|\leq1/2$, the preceding bound ensures
a nonzero first-pulse projection. Relative to that projection, a timing
error $h$ leaves the remainder $2|\sin(\rho h/2)|\leq\rho|h|$.
These statements are analytic consequences of the fixed linear model.
The following finite-time calculation tests the zero-detuning case only;
it does not numerically establish the entire detuning or timing-error family.
The common unknown modal norm cancels
from the cancellation condition; extinguishing an arbitrary pre-existing
excitation would require its complex amplitude and calibration.
Other modes, continuum waves, radiation already emitted, and nonlinear
backreaction are not canceled by this one condition. Neither zero modal
projection nor zero linear radiation implies that the charged Q-ball itself
has stopped rotating or disappeared.
